% Section 2 — Methodology.
% Demonstrates: two columns, a figure, and an equation.

\begin{frame}{Approach}

% Two columns. Handy for text beside a figure.
\begin{columns}[T]
\begin{column}{0.52\textwidth}
\begin{block}{Steps}
\begin{itemize}
    \item Data collection
    \item Pre-processing
    \item Model fitting
    \item Validation
\end{itemize}
\end{block}
\end{column}

\begin{column}{0.44\textwidth}
\begin{block}{Why this way}
\begin{itemize}
    \item A reason, not a restatement
    \item The alternative you rejected
\end{itemize}
\end{block}
\end{column}
\end{columns}

\end{frame}


\begin{frame}{A figure}

% Put image files in img/. scale= or width= both work; width relative to \linewidth
% travels better between screens and projectors.
\begin{figure}
\centering
\includegraphics[width=0.45\linewidth]{img/bracu_logo.png}
\caption{Replace this with your own figure. Captions are numbered automatically.}
\end{figure}

\end{frame}


\begin{frame}{An equation}

\begin{block}{}
One displayed equation per slide is usually the limit an audience can absorb:
\begin{equation}
    \hat{y} = \sigma\!\left(\beta_0 + \sum_{i=1}^{n} \beta_i x_i\right)
\end{equation}
Define every symbol you use, in words, immediately after it.
\end{block}

\end{frame}
